% This must be in the first 5 lines to tell arXiv to use pdfLaTeX, which is strongly recommended.
\pdfoutput=1
% In particular, the hyperref package requires pdfLaTeX in order to break URLs across lines.

\documentclass[11pt]{article}

% Remove the "review" option to generate the final version.
\usepackage[]{acl}

% Standard package includes
\usepackage{times}
\usepackage{latexsym}

% For proper rendering and hyphenation of words containing Latin characters (including in bib files)
\usepackage[T1]{fontenc}
% For Vietnamese characters
% \usepackage[T5]{fontenc}
% See https://www.latex-project.org/help/documentation/encguide.pdf for other character sets

% This assumes your files are encoded as UTF8
\usepackage[utf8]{inputenc}

% This is not strictly necessary, and may be commented out,
% but it will improve the layout of the manuscript,
% and will typically save some space.
\usepackage{microtype}

% If the title and author information does not fit in the area allocated, uncomment the following
%
%\setlength\titlebox{<dim>}
%
% and set <dim> to something 5cm or larger.

\title{Project Proposal (Put your own title here)}

% Author information can be set in various styles:
% For several authors from the same institution:
% \author{Author 1 \and ... \and Author n \\
%         Address line \\ ... \\ Address line}
% if the names do not fit well on one line use
%         Author 1 \\ {\bf Author 2} \\ ... \\ {\bf Author n} \\
% For authors from different institutions:
% \author{Author 1 \\ Address line \\  ... \\ Address line
%         \And  ... \And
%         Author n \\ Address line \\ ... \\ Address line}
% To start a seperate ``row'' of authors use \AND, as in
% \author{Author 1 \\ Address line \\  ... \\ Address line
%         \AND
%         Author 2 \\ Address line \\ ... \\ Address line \And
%         Author 3 \\ Address line \\ ... \\ Address line}

\author{First Author \\
  Affiliation / Address line 1 \\
  Affiliation / Address line 2 \\
  Affiliation / Address line 3 \\
  \texttt{email@domain} \\\And
  Second Author \\
  Affiliation / Address line 1 \\
  Affiliation / Address line 2 \\
  Affiliation / Address line 3 \\
  \texttt{email@domain} \\}

\begin{document}
\maketitle
\begin{abstract}
Abstract is optional. The proposal has to be 4 pages excluding references.
\end{abstract}

\section{Introduction}

Background and motivation of your project.

\section{Problem Definition and Method Design}

\subsection{Problem Definition}
Problem Definition

\subsection{Method Design}
Describe the methods you plan to implement to solve the above problem. E.g., which models, how to conduct training / inference.

If you plan to use any pre-trained models, specify what are they and what are their parameter sizes.


\section{Related Work}

Justified the novelty of your work based on at least three papers published in top NLP, ML or CV venues (including only EMNLP, ACL, NAACL, EACL, ICLR, ICML, NeurIPS, CVPR, ICCV, ECCV main conference or ``Findings'' papers, \textbf{not including any workshop papers}; or TACL, JMLR, TMLR, TPAMI) in the past 3 years.
Provide citations to these works, e.g. \cite{zhou-etal-2023-continual}.

\section{Evaluation Plan and Resources}

Specify the dataset, the evaluation metrics, and the expected computing resources. Justify that you have the computing resource to complete the project in the course time.

\subsection{Datasets}

Describe the dataset(s) you plan to use for your experiments, including their content and size.

\subsection{Evaluation Protocol}

Metrics, splits of training and testing. Groups of test data (if any).

\subsection{Estimated Computing Resources}

Justify that you have the computing resource to complete the project in the course time. If you plan to use a pretrained model, justify if the resource you have can support that.





% Entries for the entire Anthology, followed by custom entries
\bibliography{anthology,custom}
\bibliographystyle{acl_natbib}

%\appendix



\end{document}
